\section{ПРОГРАММНАЯ РЕАЛИЗАЦИЯ}
В данном разделе будет описан вычислительный эксперимент, 
реализация графов и обработка результатов экспериментов.
\subsection{Реализация графов и алгоритмов на них}

Так как для экспериментальной проверки свойств 
необходимы графы с большим количеством вершин,
готовые реализации графов, такие как
встроенные средства работы с графами в системе Wolfram Mathematica и 
библиотека языка программирования Python NetworkX,
не подходят из-за низкой производительности.
Поэтому графы были реализованы на языке программирования C++.

Графы были реализованы с использованием объектно-ориентированного подхода. Для этого был создан абстрактный класс <<BaseGraph>>, 
который имеет следующие методы:
\begin{enumerate}
    \item get\_n -- получение количества вершин;
    \item get\_n\_edges -- получение количества ребер;
    \item has\_edge -- проверка существования ребра между двумя вершинами;
    \item add\_edge -- добавление ребра между двумя вершинами;
    \item deg -- степень вершины;
    \item get\_neighbours -- получение списка вершин, смежных данной вершине;
\end{enumerate}
Данный абстрактный класс является родительским 
для класса <<Graph>>, который реализует неориентированный невзвешенный граф, представленный списками смежности. Использование
абстрактного класса позволит при необходимости 
заменить представление графа без изменения алгоритмов 
и вычислительного эксперимента.

Алгоритмы были реализованы с помощью функций, которые принимают 
на вход константную ссылку на граф и возвращают числовое значение. Были реализованреализованы следующие алгоритмы: поиск наибольшей компоненты связности, 
поиск количества компонент связности, поиск наибольшего пути 
в графе, найденном с помощью поиска в глубину, подсчет доли
деревьев среди всех компонент связности.
\subsection{Вычислительный эксперимент}

Для численного анализа модели $G(n,p)$ нужно 
много раз изучить, как изменяется граф при изменении $p$ 
в заданном диапазоне. Рассмотрим, как проводится вычислительный эксперимент:
\begin{enumerate}
        \item генерируется последовательность $U$ из 
            $ N = \frac{n(n-1)}{2}$ псевдослучайных чисел, равномерно распределенных на $[0;1]$;
        \item $p$ дискретно меняется в диапазоне  $[p_{\min}; p_{\max}]$;
        \item создается граф, в котором ребро под 
            номером $i = 1 \dots N$ существует, если
             $U_{i} < p$;
        \item для построенного графа вычисляются требуемые характеристики;
        \item результаты сохраняются в файл;
\end{enumerate}

В результате многократного повторения данного эксперимента 
получится файл, в котором каждая строка -- 
это результат эксперимента.

Такой подход к проведению вычислительного эксперимента позволит 
проследить эволюцию одного графа при изменении $p$.

Для программной реализации данного вычислительного эксперимента был создан класс <<RandomVariation>>. При инициализации получает 
на вход количество вершин графа, количество экспериментов, 
cписок границ диапазонов, название файла, в который будут записаны результаты, и функцию для вычисления требуемых характеристик графа. Для проведения эксперимента нужно использовать два метода:
<<variation>> для создания графов и вычислений характеристик и
<<save\_results>> для сохранения результатов в файл.

\subsection{Анализ и визуализация результатов экспериментов}

Для анализа и визуализации полученных данных использовался 
язык программирования Python и библиотеки 
numpy, scipy и matplotlib. Был реализован класс <<DataLoader>> 
для чтения данных из файлов. Классы для визуализации и анализа 
наследуются от него. Например, класс <<Plotter>> для визуализации 
изменения размера наибольшей компоненты связности, реализует
метод <<plot>> для создания графика с помощью библиотеки matplotlib.
